\documentclass[11pt]{article}
\usepackage[margin=1.1in]{geometry}
\usepackage{amsmath, amssymb, amsthm, mathtools}
\usepackage{enumitem}
\usepackage[colorlinks=true, linkcolor=blue, citecolor=blue]{hyperref}

\newtheorem{theorem}{Theorem}
\newtheorem{condition}{Condition}
\theoremstyle{definition}
\newtheorem{remark}{Remark}

\newcommand{\R}{\mathbb{R}}
\newcommand{\E}{\mathbb{E}}
\newcommand{\cA}{\mathcal{A}}
\newcommand{\cB}{\mathcal{B}}
\newcommand{\cN}{\mathcal{N}}
\newcommand{\ip}[2]{\left\langle #1, #2 \right\rangle}
\DeclareMathOperator*{\argmax}{arg\,max}
\DeclareMathOperator{\supp}{supp}
\DeclareMathOperator{\conv}{conv}

\title{What is actually proved, and what makes it work\\
  \large A self-contained guide to the convergence result}
\author{\texttt{cont\_tinkering} / \texttt{theory2}}
\date{August 2026}

\begin{document}
\maketitle

\begin{abstract}
\noindent
This note states the convergence result for mixture policies trained by magnetic
mirror descent on a continuous-action zero-sum game, and then catalogues
everything that helps it and everything that breaks it. It assumes familiarity
with finite zero-sum games and Nash equilibria, and assumes nothing about
continuous action spaces or gradient-based optimisation. It is self-contained:
every object and every hypothesis is defined here.

\medskip
\noindent
The result in one sentence:

\medskip
\noindent\fbox{\parbox{0.9\linewidth}{Under four conditions on the game, plus a
good enough initialisation, plus an exact gradient oracle, the algorithm
converges, and its limit is $\varepsilon$-Nash where $\varepsilon$ is pure
\emph{representation} bias --- the error of describing a point mass by a narrow
Gaussian --- which the algorithm drives to zero on its own at rate
$\Theta(1/t)$.}}
\end{abstract}

\tableofcontents

\part{The setting and the statement}

\section{From finite games to continuous ones}

\subsection*{What changes, and what does not}

In a finite zero-sum game with payoff matrix $A$, a mixed strategy is a point
$w$ in a simplex $\Delta_K$, the payoff is the bilinear form $w^\top A v$, and
learning dynamics move a point around inside a convex set. Everything is convex,
equilibria form a convex set, and no-regret dynamics find one from any starting
point with no assumptions on $A$ whatsoever.

Now let the action sets be continuous. Player $0$ maximises and player $1$
minimises a bounded measurable payoff
\[
  u : \cA \times \cB \to \R,
\]
where $\cA \subset \R^{n}$ and $\cB \subset \R^{m}$ are compact --- for
concreteness, intervals or boxes. A mixed strategy is now a Borel probability
measure $\mu$ on $\cA$, and the expected payoff is
$U(\mu,\nu) = \iint u \, d\mu \, d\nu$.

Two things are unchanged and one thing is destroyed.

\emph{Unchanged.} The payoff is still bilinear in the pair of strategies, the
game is still zero-sum, and an equilibrium still exists (Glicksberg's theorem, on
compact sets with continuous payoff). The strategy space is still convex --- it
is a set of measures.

\emph{Destroyed.} You cannot store a measure. Any implementable algorithm
parameterises a restricted family of measures, and the moment you do that, the
parameterisation is not convex even though the set of measures is. This is the
entire source of difficulty, and it is worth being precise about why.

\subsection*{Where an equilibrium lives}

Throughout, fix a Nash equilibrium $(\mu^\star,\nu^\star)$ and assume it is
supported on \emph{finitely many points}:
\[
  \supp \mu^\star = \{p_1,\dots,p_{M_0}\} \subset \cA,
  \qquad
  \supp \nu^\star = \{r_1,\dots,r_{M_1}\} \subset \cB,
\]
with positive weights. Call the $p_j$ the \emph{atoms}. This is not automatic ---
continuous games can have equilibria with continuous support --- but it is the
regime the algorithm is built for, and it is the common case in the games of
interest (a finite number of ``good moves'' with the right relative frequencies).

Define the \emph{effective landscape} of player $0$ against an opponent policy
$\nu$:
\begin{equation}\label{eq:landscape}
  Q_\nu(a) \;=\; \int_\cB u(a,b)\,d\nu(b)
  \qquad\text{---\; the value of playing the pure action $a$ against $\nu$.}
\end{equation}
A best response to $\nu$ is a maximiser of $Q_\nu$. The following elementary fact
is used constantly and is worth stating on its own, because it is what ties the
geometry to the game theory:

\begin{quote}
\textbf{Support characterisation.} Every atom of $\mu^\star$ maximises
$Q_{\nu^\star}$, and all atoms attain the same value. Equivalently: at
equilibrium, the atoms sit exactly on the global peaks of the landscape induced
by the opponent's equilibrium strategy, and they are tied.
\end{quote}

\noindent
This is a theorem about equilibria, not an assumption about the algorithm. It
means the target of the whole exercise can be described geometrically: \emph{find
the peaks of $Q_{\nu^\star}$, and put the right mass on each}. That is the
factorisation the entire theory runs on.

\subsection*{How far from equilibrium is a strategy pair?}

The scalar being bounded is \emph{exploitability}:
\begin{equation}\label{eq:expl}
  \varepsilon(\mu,\nu)
  \;=\;
  \Big[\max_{\mu'} U(\mu',\nu) - U(\mu,\nu)\Big]
  \;+\;
  \Big[U(\mu,\nu) - \min_{\nu'} U(\mu,\nu')\Big]
  \;=\;
  \max_{a\in\cA} Q_\nu(a) \;-\; \min_{b\in\cB} W_\mu(b),
\end{equation}
where $W_\mu(b) = \int u(a,b)\,d\mu(a)$ is the opponent's landscape and we have
used that a best response can be taken to be a pure action. A pair with
$\varepsilon = 0$ is a Nash equilibrium; a pair with $\varepsilon \le \epsilon$
is an $\epsilon$-Nash equilibrium. Every bound in this note is a bound on
\eqref{eq:expl}.

\section{The algorithm}

\subsection*{The policy}

Each player represents its strategy as a $K$-component \emph{mixture}. For player
$0$:
\begin{equation}\label{eq:policy}
  \mu_\theta \;=\; \sum_{k=1}^{K} w_k \, \cN\!\big(m_k, \sigma_k^2 I\big),
  \qquad w \in \Delta_K ,
\end{equation}
with three groups of parameters:
\begin{itemize}[leftmargin=1.4em, itemsep=2pt]
\item the \emph{categorical head} $w$ --- how much mass on each component;
\item the \emph{means} $m_1,\dots,m_K \in \cA$ --- where the components sit;
\item the \emph{widths} $\sigma_1,\dots,\sigma_K > 0$ --- how spread out each is,
  floored at some $s^{\min} > 0$.
\end{itemize}
The widths are genuine parameters, updated by the same objective as everything
else; they are not a schedule and not a hyperparameter. This will matter enormously.

Read \eqref{eq:policy} against the finite case: $w$ is the familiar mixed
strategy, and $(m,\sigma)$ answer a question that does not exist in a finite game
--- \emph{which pure actions am I mixing over?} In a finite game the support is
given by the game. Here it must be discovered, and discovering it is a nonconvex
problem in $\cA$.

\subsection*{The update: magnetic mirror descent}

Both players run magnetic mirror descent (MMD). On the categorical head this is
mirror descent (multiplicative weights) with entropy regularisation at
temperature $\tau$, plus a \emph{magnet}: a proximal term pulling toward a
snapshot of the strategy taken every $T$ steps, which is periodically refreshed
to the current iterate. On the means and widths it is preconditioned gradient
ascent, optionally with an analogous Gaussian magnet at temperature
$\tau^{\mathrm g}$.

Two facts about the magnet, both used later:
\begin{itemize}[leftmargin=1.4em, itemsep=2pt]
\item On a fixed \emph{matrix} game the refreshed magnet is a proximal-point
  method whose fixed points are exactly the Nash equilibria, for \emph{every}
  $\tau > 0$. So the entropy temperature biases the \emph{rate} and not the
  \emph{limit}. This is not obvious --- one expects entropy regularisation to
  leave an $O(\tau)$ residual --- and it is the refresh that removes it.
\item On the \emph{means}, the Gaussian magnet does something quite different,
  and less pleasant. See \S\ref{sec:brake}.
\end{itemize}

\subsection*{The one computation that makes the analysis possible}

Write $\bar A_{kl}$ for the expected payoff when player $0$ samples from its
component $k$ and player $1$ from its component $l$. Then, by Fubini alone, with
no assumption on $u$ and none on the component shapes,
\begin{equation}\label{eq:decomp}
  U(\mu_\theta, \nu_\theta) \;=\; w^\top \bar A\, v ,
\end{equation}
and the gradients the two parameter groups receive are:
\begin{enumerate}[leftmargin=1.8em, itemsep=2pt]
\item[(i)] the categorical head sees $\bar A v$ --- so it is running
  \emph{verbatim} tabular MMD on the finite matrix game $\bar A$;
\item[(ii)] the mean $m_k$ sees $w_k \nabla_{m_k}\, \E_{a \sim \cN(m_k,\sigma_k^2)}
  Q_\nu(a) = w_k \nabla (Q_\nu)_{\sigma_k}(m_k)$ --- gradient ascent on the
  landscape \emph{blurred} at that component's own width, where
  \begin{equation}\label{eq:blur}
    \Lambda_s \;:=\; \Lambda * \cN(0, s^2 I)
  \end{equation}
  denotes convolution with a Gaussian of width $s$.
\end{enumerate}

Clause (i) is the important one. The categorical head is not \emph{approximately}
playing a matrix game, or playing one in the limit; it is playing one exactly, at
every step, wherever the components happen to sit. Clause (ii) says each mean
climbs a blurred version of the landscape, on its own, seeing nothing of the other
components.

So the problem factorises:
\begin{center}
\begin{tabular}{@{}lll@{}}
\hline
question & object & difficulty \\
\hline
how much mass on each atom? & a finite matrix game & \emph{solved, unconditionally} \\
where are the atoms? & nonconvex ascent in $\cA$ & \emph{everything hard} \\
\hline
\end{tabular}
\end{center}

\noindent
Every hypothesis in the theory belongs to the second row.

\section{Why this is not simply ``run MMD''}

A game theorist's first instinct is that MMD is a no-regret method on a convex
set and should therefore just work. Four obstructions say otherwise, and it is
worth seeing all four before the conditions arrive, because each condition is a
response to one of them.

\paragraph{1. No single mean can converge to an equilibrium.}
An equilibrium is a \emph{measure} over $M_0 > 1$ atoms. A mean $m_k$ is a
\emph{point} of $\cA$, and no point of $\cA$ is an equilibrium. So there is no
hope of a condition of the form ``each mean is attracted to the equilibrium''.
The only shape a condition can take is: \emph{each mean is attracted to one
atom} --- and then mixing over the atoms is not a geometric phenomenon at all, it
is the matrix game of \eqref{eq:decomp}. This forces the analysis to partition
$\cA$ into regions, one per atom.

\paragraph{2. The landscape is not concave, and local ascent is local.}
$Q_\nu$ is whatever the game makes it. It generally has several local maxima.
Gradient ascent finds the nearest one and stays there. Nothing in the update
compares one peak's height to another's, so a component that starts near a bad
peak converges to the bad peak, forever, regardless of how much worse it is.

\paragraph{3. The landscape moves while you climb it.}
$Q_\nu$ is built from the opponent's current policy. As the opponent learns, the
landscape tilts and its peaks slide. You are not climbing a hill; you are chasing
one. Worse, this breaks the obvious repair to obstruction 2: one would like to
assume ``each mean starts in the basin of attraction of its atom'', but basins of
a \emph{moving} landscape are not forward-invariant. A component correctly placed
at time $t$ can find itself in the wrong basin at time $t+1$ because the opponent
moved, through no fault of its own and with no update applied to it.

\paragraph{4. The policy class does not contain its target.}
The target is atomic; a Gaussian of width $\sigma \ge s^{\min} > 0$ is not an
atom. Exact convergence to $\mu^\star$ is therefore impossible in the
parameterisation actually implemented --- not for algorithmic reasons but for
representational ones. Any honest theorem must produce an $\varepsilon$, and must
say where it comes from.

\section{The four conditions}\label{sec:conditions}

\subsection*{(C): one attractor per cell}

This is the central hypothesis and the one that excludes games.

\begin{condition}[Attracting cell structure]\label{cond:C}
There are open sets $\cA_1,\dots,\cA_{M} \subset \cA$ with disjoint closures
whose union has full measure, and points $\xi_j \in \cA_j$, such that each
$\cA_j$ is \emph{star-shaped about} $\xi_j$ (for every $a \in \cA_j$ the whole
segment from $\xi_j$ to $a$ stays in $\cA_j$) and, for every $j$ and every
$a \in \cA_j$,
\begin{equation}\label{eq:star}
  \ip{\nabla \Lambda(a)}{\xi_j - a}
  \;\ge\;
  \rho\big(\|\xi_j - a\|\big)\;\|\xi_j - a\| ,
\end{equation}
for some continuous $\rho$ with $\rho(0) = 0$ and $\rho(r) > 0$ for $r > 0$.
\end{condition}

\noindent
In words, \eqref{eq:star} says: \emph{from anywhere in the cell, walking uphill
takes you at least slightly towards that cell's target.} The gradient may be
arbitrarily small, may point mostly sideways, and the landscape may have flat
stretches and inflection points. What it forbids is exactly two things --- a
second attractor inside the cell, and the gradient pointing \emph{away} from
$\xi_j$ anywhere inside it.

Three clarifications, each of which is a common misreading.

\begin{itemize}[leftmargin=1.4em, itemsep=3pt]
\item \textbf{It is far weaker than concavity.} A Gaussian bump
  $h\,e^{-x^2/2w^2}$ is concave only within $|x| \le w$, yet it satisfies
  \eqref{eq:star} on its \emph{entire} cell, because its gradient points uphill
  everywhere. Concavity compares values at every pair of points;
  \eqref{eq:star} merely signs one inner product per point.
\item \textbf{$\xi_j$ is posited, not defined.} To claim a game qualifies you must
  \emph{exhibit} the cells and the targets. That $\xi_j$ is then the unique
  critical point and the unique maximiser of $\Lambda$ over $\overline{\cA_j}$ is
  a \emph{consequence} of \eqref{eq:star}, not an extra assumption. This is what
  makes (C) checkable by counting peaks rather than by solving for critical
  points.
\item \textbf{The landscape here is the blurred one, and the condition is
  quantified.} What the full argument requires is \eqref{eq:star} for
  $\Lambda = (Q_\nu)_s$, for \emph{every} width $s \in [s^{\min}, s^{\max}]$ and
  \emph{every} opponent policy $\nu$ the run can visit --- \textbf{with the same
  cells throughout}. That last clause is not a technicality. It is the answer to
  obstruction 3: a cell is morally the intersection of the basins over the whole
  family of landscapes, the region belonging to atom $j$ \emph{no matter what the
  opponent does}. That is what supplies a single measure of progress valid for
  all time.
\end{itemize}

\subsection*{(L): the peaks are nondegenerate}

\begin{condition}[Linear modulus]\label{cond:L}
$\rho(r) \ge \alpha r$ for some $\alpha > 0$, on the whole cell.
\end{condition}

\noindent
On a bounded cell this is equivalent to requiring the peaks to be
\emph{nondegenerate}, uniformly --- no flat-topped maxima. It is needed because
the algorithm takes discrete steps of a fixed size, not infinitesimal ones; see
\S\ref{sec:whyL}. Note it does \emph{not} require $\alpha$ to be large.

\subsection*{(M): the peaks found are the atoms}

\begin{condition}[Peak--atom identification]\label{cond:M}
At the equilibrium opponent $\nu^\star$ and the terminal width, the targets
$\xi_1,\dots,\xi_M$ of the cells are exactly the equilibrium atoms
$p_1,\dots,p_{M_0}$. In particular $M = M_0 = |\supp \mu^\star|$.
\end{condition}

\noindent
Half of this is free. By the support characterisation, the atoms are
\emph{automatically} maxima of $Q_{\nu^\star}$ --- that is a theorem. What (M)
asserts is the \emph{converse pairing}: that the attractors the cell structure
provides are those maxima \emph{and no others}. Combined with (C)'s insistence
that there are no other attractors, the two say together:
\begin{quote}
\emph{the landscape's peaks and the equilibrium's atoms are the same finite set.}
\end{quote}
This is where the geometry becomes game theory, and it is the step at which a
statement about hill-climbing becomes a statement about Nash equilibria.

\subsection*{(G): the mutual disturbance damps}

Four groups of parameters move at once: both players' means and both players'
weights. Each disturbs the others --- your weights change your opponent's
landscape, your opponent's means move your targets. Let $x_i(t)$ be the error of
block $i$ (distance from its own target), and suppose each block satisfies an
inequality of the form
\begin{equation}\label{eq:blocks}
  x_i(t+1) \;\le\; \underbrace{(1 - a_i)\,x_i(t)}_{\text{self-repair}}
  \;+\; \underbrace{\textstyle\sum_{j \ne i} b_{ij}\,x_j(t)}_{\text{disturbance received}}
  \;+\; \underbrace{d_i(t)}_{\text{external}} .
\end{equation}

\begin{condition}[Small gain]\label{cond:G}
The $4 \times 4$ matrix $\hat\Gamma_{ij} := b_{ij}/a_i$ --- disturbance received
per unit of the receiver's own repair rate --- has spectral radius below one.
\end{condition}

\noindent
This is the standard small-gain principle: an interconnection of contracting
subsystems is stable if no loop amplifies. If $\rho(\hat\Gamma) < 1$ there is a
weighted maximum $V = \max_i x_i / \pi_i$ that contracts, and $x(t) \to 0$
\emph{provided the external inputs $d_i(t)$ vanish}. Keep that proviso in mind; it
is doing real work in \S\ref{sec:anneal}.

The two-block version, with only the means moving, is easier to state and worth
carrying as intuition. Let $g_0$ be the sensitivity of player $0$'s attractor to a
displacement of player $1$'s means, and $g_1$ symmetrically. The chase damps when
\begin{equation}\label{eq:sg2}
  g_0 \, g_1 \;<\; 1 ,
  \qquad\text{where}\qquad
  g_0 \;\le\; \frac{\text{coupling strength}}{\text{curvature } \alpha_0} .
\end{equation}
\emph{Curvature resists; coupling drives.} A sharply peaked landscape holds its
peak in place when the opponent tilts it; a shallow one lets the peak slide far.

This is the one genuinely game-theoretic condition in the development: it has no
counterpart in single-agent optimisation, because there is no opponent and hence
no round trip.

\section{The result}

\begin{theorem}[Convergence]\label{thm:main}
Assume (C) and (L) uniformly over the admissible width range and over every
opponent policy the run can visit, with the same cells; assume (M); assume (G).
Assume in addition:
\begin{enumerate}[leftmargin=1.8em, itemsep=2pt]
\item[(a)] \emph{covering initialisation}: every cell contains at least one
  component mean, and that mean lies within the cell's \emph{invariance radius}
  $R_j := \mathrm{dist}(\xi_j, \partial \cA_j)$;
\item[(b)] \emph{admissible hyperparameters}: step size $\eta < \alpha/L^2$
  where $L$ is the Lipschitz constant of $\nabla\Lambda$, and entropy temperature
  $\tau$ large enough to satisfy (G);
\item[(c)] \emph{exact gradients}.
\end{enumerate}
Then every mean converges to its cell's target, every width falls to the floor
$s^{\min}$ in finitely many steps, the component matrix $\bar A$ converges to the
smoothed atom matrix, and the categorical heads converge to a Nash equilibrium of
it. The means are within $O\big((s^{\min})^2\big)$ of the atoms, and the induced
policy pair is $\varepsilon$-Nash with
\begin{equation}\label{eq:eps}
  \varepsilon \;=\; O\!\big((s_0^{\min})^2 + (s_1^{\min})^2\big) .
\end{equation}
Released from the floor, the widths obey $\dot\sigma = -c\,\sigma^3$ in closed
form, so $\sigma(t) = \Theta(t^{-1/2})$ and $\varepsilon(t) = \Theta(1/t)$.
\end{theorem}

\subsection*{Unpacking the one-liner}

\paragraph{``Under conditions on the game''.}
(C), (L), (M) are properties of the payoff $u$. (G) is a $4\times4$ eigenvalue
computation whose entries are curvatures and coupling norms of $u$ --- crucially,
it contains \emph{no step size and no width}, so it is checkable before training
rather than along the trajectory. Together they define a class of games. A game
either is in the class or is not, and no amount of tuning moves it.

\paragraph{``Plus a good enough initialisation''.}
Condition (a) is not a property of the game and cannot be achieved by the
algorithm; it is assumed. Note it is strictly stronger than ``each cell contains
a component''. Being merely \emph{inside} a cell is not enough: the argument
guarantees that the \emph{distance to the target} does not grow, which confines
the trajectory to the cell only if the ball of that radius fits inside the cell.
Otherwise a component can exit through a nearby piece of the cell boundary while
moving strictly closer to its target, land in a neighbouring cell, and converge to
the wrong attractor.

\paragraph{``Plus an exact gradient oracle''.}
The largest gap between the theorem and any implementation. See \S\ref{sec:gaps}.

\paragraph{``Converges''.}
All four blocks reach a limit: two sets of means, two sets of weights (with the
widths reaching the floor in finite time along the way). This is the clause that
(C), (L), (G) and the initialisation buy. They buy nothing else --- none of their
constants appears in \eqref{eq:eps}.

\paragraph{``$\varepsilon$-Nash, where $\varepsilon$ is pure representation bias''.}
This is the clause that surprises people, so it deserves the most space. The
$\varepsilon$ is not slack in the analysis and not a residual from the entropy
regularisation. It is exactly obstruction 4 of \S3, and one can see this by
changing the component family and watching it vanish:

\begin{quote}
Run the same algorithm with \textbf{point masses} $\delta_{m_k}$ instead of
Gaussians. Then $\bar A_{kl} = u(m_k, n_l)$ \emph{exactly}, so with the means at
the atoms $\bar A$ \emph{is} the atom matrix, and $\varepsilon = 0$ exactly ---
the equilibrium is a fixed point of the algorithm, not a limit of one.
\end{quote}

\noindent
So $\varepsilon$ appears for one reason only: a Gaussian of width $s^{\min} > 0$
is not a point mass, so the policy class does not contain its target. Concretely,
\eqref{eq:eps} decomposes into two effects of the same width, and it is worth
separating them:
\begin{itemize}[leftmargin=1.4em, itemsep=2pt]
\item \emph{the target is displaced}: the peak of the \emph{blurred} landscape
  $\Lambda_s$ is not the peak of $\Lambda$; it sits $O(s^2)$ away. This is the
  ``means within $O(s^2)$ of the atoms'' clause.
\item \emph{the value is deficient}: playing a spread-out component around a peak
  earns less than playing the peak itself, by the \emph{blur deficit}
  $\max_\cA \Lambda - \max_\cA \Lambda_{s^{\min}}$, which is $O(s^2)$ whenever
  the landscape is twice differentiable at its peaks.
\end{itemize}

\paragraph{``Which the algorithm drives to zero on its own at $\Theta(1/t)$''.}
And this is the reason the trade is worth making. Nothing schedules the width
down; the width has its own gradient, and \S\ref{sec:anneal} explains why that
gradient always points down near a peak. Remove the floor and the bias
extinguishes itself as a \emph{rate} rather than requiring a tuned schedule. The
constant is available in closed form, and the predicted decay law has been
verified numerically to four digits.

\subsection*{Which part of the result each hypothesis buys}

This is the single most useful table in the note.

\begin{center}
\begin{tabular}{@{}p{0.26\textwidth}p{0.36\textwidth}p{0.30\textwidth}@{}}
\hline
hypothesis & buys & appears in the final bound? \\
\hline
matrix-game factorisation & the form of $\varepsilon$ & \textbf{yes} --- it is the formula \\
(C) one attractor per cell & convergence happens & no \\
(L) nondegenerate peaks & convergence in discrete time; a rate & no \\
(M) peaks are atoms & the limit is an \emph{equilibrium} & no \\
(G) small gain & the four blocks do not destabilise each other & no \\
covering initialisation & the right cells are reached at all & no \\
Gaussian components & self-annealing, and the $\Theta(1/t)$ law & \textbf{yes} --- it is the $O(s^2)$ \\
\hline
\end{tabular}
\end{center}

\noindent
Read the right-hand column: \emph{almost nothing that makes the algorithm work
affects how good the answer is.} Tightening $\alpha$, or $g_0g_1$, or $\tau$ makes
you converge faster or brings more games into the class. It does not land you
closer to Nash. Only the component family and the width floor do that.

\part{What makes it work}

Six mechanisms. Each is a specific structural fact, and each is load-bearing in
the sense that the argument fails without it.

\section{Mixing is free}

The factorisation \eqref{eq:decomp} is exact and assumption-free: whatever the
components are and wherever they sit, the categorical heads play a genuine finite
matrix game, and finite zero-sum matrix games are solved by MMD globally, from
any initialisation, at a linear rate, with no assumptions on the matrix.

This halves the problem in a very strong sense. Run the argument with the means
in completely the wrong places and it still goes through --- the weights converge,
beautifully and globally, to the Nash equilibrium of a matrix that describes the
wrong components. That is why every hypothesis in \S\ref{sec:conditions} is about
locations.

It also explains the shape of the final bound. At the limit, the heads solve
$\bar A$, and a Nash equilibrium of $\bar A$ \emph{equalises} the components in
its support (complementary slackness): every supported component attains the same
value. So the mixture's payoff equals the best single component's payoff, and the
only remaining gap is between that and the best achievable value at a
\emph{point} --- which is precisely the blur deficit. The $\varepsilon$ formula is
not derived from any dynamic argument; it falls out of complementary slackness.

\section{Steps shrink as components arrive}\label{sec:vanish}

This is the most quietly important fact in the whole development. Condition (C)
implies, with no further assumption, that
\begin{equation}\label{eq:dominated}
  \|\nabla\Lambda(a)\| \;\le\; L\,\|a - \xi_j\|
  \qquad\text{for every } a \in \cA_j ,
\end{equation}
where $L$ is the Lipschitz constant of the gradient. (Reason: (C) forces $\xi_j$
to be a critical point, and a Lipschitz gradient vanishing at $\xi_j$ is bounded
by $L$ times the distance to it.)

In words: \emph{a component's step shrinks as it arrives}. Why this matters is
clearest in the four-block system \eqref{eq:blocks}. The disturbance each block
injects into its neighbours \emph{is} its own step. If steps were merely of size
$O(\eta)$ --- as one would expect from a fixed-step method --- then $b_{ij}$ would
have a floor and the four blocks would settle into a permanent mutual jitter of
size $O(\eta)$ rather than converging. Because of \eqref{eq:dominated} the
disturbance is proportional to the emitting block's \emph{error}, so it vanishes
as that block arrives, and the interconnection has a genuine fixed point at zero
error.

Every later mechanism relies on this one.

\section{The blur anneals itself}\label{sec:anneal}

Here the Gaussian is bought, and it is bought for two specific properties that no
other component family supplies.

\paragraph{The width is a parameter, not a hyperparameter.} It is a coordinate of
the policy, and the same objective that moves the means moves it. Contrast an
algorithm whose exploration bandwidth is exogenous: there is something to tune,
not something that moves.

\paragraph{The Gaussian is the heat kernel.} Blurring at width $s$ is exactly
running the heat equation for time $s^2/2$:
\begin{equation}\label{eq:heat}
  \frac{\partial}{\partial (s^2/2)}\, \Lambda_s \;=\; \Delta \Lambda_s .
\end{equation}
At an interior maximum the Laplacian is non-positive, and strictly negative if
the maximum is nondegenerate. So: \emph{blurring a peak lowers it}. Therefore the
width update, which is ascent on value, is strictly negative near a peak, and the
width falls. Nobody schedules this. It is a consequence of the Gaussian being a
semigroup for the Laplacian, and a family without that identity has no reason for
its spread to fall at a peak rather than rise.

Two upgrades turn \eqref{eq:heat} from a curiosity into a usable mechanism.

\emph{From ``at the peak'' to ``near the peak''.} Nondegeneracy makes the
Laplacian bounded away from zero throughout a ball around the peak, not just at
it. This matters because transport and annealing happen \emph{simultaneously} ---
one never gets to apply \eqref{eq:heat} at an exactly-arrived component.

\emph{Finite arrival time.} Consequently the width reaches the floor in
\emph{finitely many steps}, after which it is constant and contributes
\emph{exactly zero} disturbance. Recall the proviso attached to (G): the
conclusion $x(t) \to 0$ holds only if the external inputs $d_i(t)$ vanish. The
widths are those inputs, and it is finite arrival --- not asymptotic decay --- that
makes them vanish. This is the difference between concluding \emph{convergence}
and concluding \emph{convergence to a ball}.

\paragraph{The obvious objection, and why it fails.} The mean update is
preconditioned by $\sigma^2$, so as the width collapses from $10^{-1}$ to
$10^{-3}$ the mean's progress per step falls by four orders of magnitude. A
component that has not arrived before the width collapses ought to freeze where it
stands. But the \emph{disturbance} the collapse injects falls by \emph{six}
orders: one factor $s^2$ because the target moves that slowly per unit of
blur-time, and one more from the preconditioner in the width update itself. Only
the ratio matters, and the ratio improves as the widths fall. Transport is not
throttled relative to what it is chasing; both slow down, and the disturbance
slows faster.

\section{The preconditioner cancels, so the stability condition is about the game}

The same cancellation has a second, larger consequence. In the four-block gain
matrix $\hat\Gamma_{ij} = b_{ij}/a_i$, the factor $\sigma^2$ appears once in a
numerator and once in a denominator around every cycle of the interconnection. It
therefore cancels, and (G) contains \emph{no width and no step size}.

That is what makes (G) a statement about the game --- a finite eigenvalue
computation from curvatures and coupling norms, done before training --- rather
than a statement about a trajectory, which would be unverifiable.

\section{The entropy temperature is free}

Two of the cycles in $\hat\Gamma$ are binding. One is the mean-versus-mean cycle,
which is \eqref{eq:sg2}: curvature against coupling, a property of the game. The
other is a mean-versus-weight cycle, and it carries a factor $1/\tau$: the larger
the entropy temperature, the stiffer the weight block's target is against the
means moving underneath it, and the less it is disturbed.

So that cycle can always be closed by turning $\tau$ up. Ordinarily this would be
a bad trade --- more regularisation, more bias. It is not, because of the magnet:
at a fixed point of the magnet refresh the regularisation term vanishes
identically, so $\tau$ moves the rate and not the limit. A knob that buys
stability at no asymptotic cost is rare, and this one is genuinely free.

\section{The width floor buys breadth, not just stability}

The floor $s^{\min}$ looks like an implementation detail that costs accuracy. It
is the opposite on both counts.

\emph{It costs no accuracy}, because the floor is not what drives the bias down
--- the width dynamics do, and removing the floor recovers exactness as a rate
($\varepsilon = \Theta(1/t)$).

\emph{It buys breadth}, in two ways:
\begin{itemize}[leftmargin=1.4em, itemsep=2pt]
\item Condition (C) is imposed from $s^{\min}$ upward rather than from $0$. So a
  landscape carrying structure \emph{finer} than the floor --- ripple,
  discretisation artefacts, a Monte-Carlo estimate --- can have very many peaks at
  width $0$ and exactly $M$ at the floor. Such a game is inside the class with a
  floor and outside it without.
\item It is what allows the payoff to be merely \emph{bounded and measurable}
  rather than twice continuously differentiable. Everything the algorithm ever
  evaluates is blurred at width at least $s^{\min}$, hence smooth with bounded
  derivatives no matter how irregular $u$ is.
\end{itemize}

\part{What breaks it}

\section{Fatal: a second attractor}\label{sec:fatal}

This is the failure mode the whole class is designed around, and it is worth
stating why it cannot be engineered away.

Suppose some cell contains a second nondegenerate local maximum $\hat a \ne \xi_j$
--- a ``decoy''. Then:

\begin{enumerate}[leftmargin=1.8em, itemsep=3pt]
\item[(i)] $\hat a$ is an exponentially stable fixed point of the mean dynamics,
  with an \emph{open} set of initialisations converging to it. Local ascent near
  a nondegenerate maximum is a contraction, and contraction is an open condition.
  So the trap has positive measure; a random initialisation hits it with positive
  probability.
\item[(ii)] It is \textbf{magnet-invariant}. The magnet snapshot is refreshed to
  the current iterate, so at the trap the magnet's own gradient vanishes
  identically --- and the magnet contributes a negative-definite term to the
  Jacobian, strictly \emph{deepening} the trap. Reaching for the magnet makes it
  worse.
\item[(iii)] It is \textbf{capacity-invariant}. Spare components initialised in
  its basin are captured too, so raising $K$ does not help.
\item[(iv)] It is \textbf{annealing-invariant}, and worse: an annealing schedule
  is the \emph{mechanism that delivers} components into it, since which branch a
  component follows down the scale ladder is decided at large blur, where the
  landscape's coarse structure dominates.
\end{enumerate}

\paragraph{Why ``dominated'' is not enough.} The natural weaker hypothesis is
``every spurious peak is lower than every atom''. That controls which peak is the
\emph{global} maximum --- but local ascent never compares global amplitudes. A
dominated spurious bump still captures whatever enters it. This is why the
condition has to forbid extra attractors outright rather than merely rank them,
and it is the sense in which (C) is \emph{derived} rather than chosen: \emph{any
class of games robust to initialisation must forbid spurious attractors}.

\paragraph{Why you cannot rescue it by adding a cell and a component.} A tempting
reading is: give the decoy its own cell, and use $M+1$ components. The decoy is
then its own attractor in its own cell, and (C) is satisfied. This does not work,
for a specific reason: $M$ is not free. It is pinned by (M) to
$|\supp \mu^\star|$. The re-partition satisfies (C) with $M+1$ cells but now one
target is not an equilibrium atom, so (M) fails. The violation moves; it does not
disappear. And by (iii) the extra component does not reliably go to the decoy
anyway --- capture is decided by initialisation, so you get a random split, not a
sacrificial component.

\paragraph{The test, and it is cheap.} Count the local maxima of the landscape at
the floor width and compare with the support size of the equilibrium. More peaks
than atoms, and the game is out of the class --- for good, not pending better
hyperparameters. A worked instance: a ``decoy game'' whose landscape has three
local maxima while its equilibrium has support size two. (C) asks for exactly two
cells each containing one attractor, which is impossible; some cell contains two.
There is therefore an open set of initialisations from which the method does not
converge to Nash --- which is exactly what is observed. No dominance computation
and no smoothing analysis is required: a mode count against a support size.

\paragraph{Reading it in the other direction.} There is a companion inequality
that constrains a hyperparameter rather than the game. (C) provides $M$ distinct
targets at every admissible width and (M) identifies them with the atoms, so the
atoms must be \emph{resolvable} at the floor. For a landscape of equal Gaussian
bumps of width $\bar w$, the classical bimodality threshold gives
\begin{equation}\label{eq:resolution}
  \sqrt{\bar w^2 + (s^{\min})^2} \;<\; \tfrac12 \min_{j \ne j'}\|p_j - p_{j'}\| :
\end{equation}
\emph{do not choose a floor larger than half the atom spacing.}

\section{Fatal: coupling stronger than curvature}

If \eqref{eq:sg2} fails, the chase does not damp. Each player's target runs
further than the other player can close, and the pair can cycle indefinitely ---
the continuous-action analogue of the familiar cycling of learning dynamics in
matching pennies, except that here it is the \emph{supports} that cycle, not the
mixing weights.

The condition is a genuine constraint but an easy one to satisfy in practice. For
a payoff of the form
$u(a,b) = D(a) - D(b) + c\,\ip{f(a)}{f(b)}$ --- a per-player ``self'' term plus a
coupling of strength $c$ --- with Gaussian bumps of height $h$ and width $w$ the
curvature is $\alpha \approx h/w^2$ and the condition reads
$c\,\|f'\|^2 < (h/w^2)^2$. On a concrete two-peak instance
($h=1$, $w=0.1$, so $\alpha = 100$; $f(a) = 2a-1$, $c = 1$) one gets
$g_0 = 0.02$ and $g_0 g_1 = 4\times10^{-4}$: three orders of magnitude of margin,
with trouble first expected around $c \approx 50$. Consistent with this, no
experiment at $c=1$ has shown coupling-induced instability, and on a three-peak
variant the predicted onset at $c \approx 78$ matches an observed degradation
between $c=32$ and $c=300$.

The useful summary is that the coupling cap scales as $h/w^2$: \emph{sharp peaks
tolerate strong coupling; shallow ones do not}.

\section{Fatal: flat games, where there is no gradient at all}

There is a degenerate case worth knowing because it is not rare in game theory:
a payoff with \emph{no self-term},
\[
  u(a,b) \;=\; \ip{\varphi(a)}{C\,\psi(b)} ,
\]
a pure coupling. At equilibrium the landscape is \emph{constant}:
$\nabla\Lambda \equiv 0$, and \eqref{eq:star} fails at every point. There is
nothing to ascend. This is the continuous-action version of matching pennies, and
the difficulty is exactly the familiar one --- no pure action is preferred --- but
now it manifests as the transport mechanism having nothing to work with.

Here the Gaussian magnet is not optional but necessary: its quadratic term
$-\tfrac{\tau^{\mathrm g}}{2}\|a - \bar a\|^2$ supplies the curvature the game
does not have, and (C) then holds with one cell. The price is real: the target is
the previous iterate rather than a feature of the game, so the atoms are
\emph{unidentified} and the conclusion degrades from ``the means converge to the
atoms'' to ``the induced measure converges''.

Two riders. First, the right object for (C) is therefore not the blurred landscape
but the \emph{magnet-regularised} blurred landscape --- curved and flat games are
then two extremes of one condition rather than two separate cases. Second, the
dichotomy is per \emph{direction}, not per game: a two-dimensional action can
have a well in one coordinate and matching pennies in the other, and a scalar
$\tau^{\mathrm g}$ then pays the price in both.

\section{Costly: the Gaussian magnet is a pure brake}\label{sec:brake}

On the categorical head the magnet is free (\S8). On the \emph{means and widths}
it is not, and the difference is worth stating because reaching for the magnet is
the natural response to slow transport.

With magnet temperature $\tau^{\mathrm g}$ refreshed every $T$ steps, set
\[
  \beta \;=\; \frac{1 - e^{-\eta\tau^{\mathrm g}T}}{\eta\tau^{\mathrm g}T}
  \;\in\; (0,1] .
\]
Then \emph{both} Gaussian-head coordinates --- means and widths --- are slowed by
the same factor $\beta$: the dynamics with the magnet are the dynamics without it
under the time reparameterisation $t \mapsto \beta t$.

A time reparameterisation cannot change which initialisations eventually
converge. So the consequences are stark: asymptotic basins are \emph{unchanged},
and what shrinks is finite-horizon reach. At a fixed step budget you arrive at
fewer targets, and the bias at a fixed step count is worse by $1/\beta$. Combined
with \S\ref{sec:fatal}(ii) --- the magnet makes a trap deeper --- the verdict for a
curved landscape is that the Gaussian magnet cannot help transport and can only
cost. (In a flat game it is necessary, per the previous section; those are the two
extremes.)

The factor has been checked over a $16\times$ range in the product
$\eta\tau^{\mathrm g}T$: configurations sharing that product collapse onto the
same constant, confirming $\beta$ depends on nothing else. Predicted versus
measured slowdown agrees to within $0.4\%$ across
$\eta\tau^{\mathrm g}T \in \{0.5, 2, 8\}$.

\section{Costly: weight starvation}

The implemented preconditioner on the mean update is $\sigma_k^2 w_k$ --- the
factor $w_k$ arriving from clause (ii) of \eqref{eq:decomp}. That $w_k$ is a
defect, not a feature. The exact natural gradient under the Fisher metric of the
joint law $p(k,a) = w_k \cN(a \mid m_k, \sigma_k^2)$ is $\sigma_k^2
\nabla\Lambda(m_k)$ with \emph{no} $w_k$.

The discrepancy has a concrete failure mode: \emph{a component losing weight also
loses mobility}, so it freezes wherever it happens to be, which may be a
non-atom. The analysis has to carry a floor $w_k \ge w_{\min} > 0$ as a
hypothesis to exclude this. Under the corrected preconditioner $\sigma_k^2/w_k$
that hypothesis becomes \emph{vacuous}.

This is worth flagging separately from the other entries in this part because it
is the one place where the difficulty is a fixable defect in the implementation
rather than a property of any game.

\section{Costly: a bad initialisation, permanently}

Local ascent cannot move a component from one cell to another. Which attractor a
component finds is decided at initialisation, once, and never revisited. This is
a genuine limitation of the algorithm and not an artefact of the analysis.

Two consequences:
\begin{itemize}[leftmargin=1.4em, itemsep=2pt]
\item \emph{Coverage must be assumed.} The honest response is to assume it and
  note that overcapacity ($K > M$) with a spread initialisation buys it cheaply
  --- but only \emph{inside} the class, where by \S\ref{sec:fatal} there is no
  spurious attractor for the spare components to fall into. Outside the class,
  overcapacity feeds the trap.
\item \emph{Cell membership is not enough.} As noted in \S6, the requirement is
  membership within the invariance radius $R_j$, which can be far smaller than the
  cell's extent --- a star-shaped cell can be a thin spike.
\end{itemize}

\section{Not fatal, contrary to appearance: a small $\alpha$}

Condition (L) requires $\alpha > 0$; it does \emph{not} require $\alpha$ to be
large, and typically it is not. For a Gaussian bump of height $h$ and width $w$
the modulus behaves like $\rho(r) \asymp (h/w^2)\,r\,e^{-r^2/2w^2}$, so
\[
  \alpha \;\asymp\; (h/w^2)\,e^{-\mathrm{diam}^2/2w^2}
\]
--- exponentially small, reflecting the ``dead zone'' far from the peak where the
landscape is essentially flat.

Every conclusion of Theorem \ref{thm:main} is uniform in $\alpha > 0$: only the
\emph{rate} degrades. This is a genuinely useful demotion. What in a cruder
theory would be a hypothesis --- some reachability radius beyond which the method
is not claimed to work --- becomes a rate, and it matches the observation that
borderline configurations converge at $150$k steps and not at $40$k. They were
never failing; they were slow.

The distinction that matters is therefore \emph{not} between large and small
$\alpha$, but between $\rho$ decaying at long range (harmless, and the generic
case) and $\rho$ vanishing \emph{superlinearly at the target}, which is what (L)
excludes and which a nondegenerate maximum never does.

\section{Why (L) cannot be dropped}\label{sec:whyL}

It is natural to read (L) as buying a rate --- geometric decay rather than mere
monotone approach --- and to expect that without it one still converges, slowly.
In continuous time that is right. With a \emph{fixed} discrete step it is not, and
the reason is instructive about the difference between the two settings.

One step of ascent splits into a radial gain and a tangential cost:
\[
  \underbrace{\eta\,\ip{\nabla\Lambda}{\xi_j - m}}_{\text{gain} \;\ge\; \eta\,\rho(\|d\|)\|d\|}
  \qquad\text{versus}\qquad
  \underbrace{\tfrac{\eta^2}{2}\|\nabla\Lambda\|^2}_{\text{cost} \;\le\; \tfrac{\eta^2}{2}L^2\|d\|^2} ,
  \qquad d := m - \xi_j .
\]
The cost is \emph{quadratic} in the distance, by \eqref{eq:dominated}. The gain is
quadratic only if $\rho(r)/r$ is bounded below --- which is exactly (L). For a
modulus vanishing faster than linearly, say $\rho(r) = r^{1+\delta}$, the gain
dominates only while $\|d\|^\delta > \tfrac{\eta}{2}L^2$, and inside that radius a
fixed step size certifies nothing.

So the honest division is: the continuous-time flow converges under (C) alone;
the fixed-step iteration needs (C) and (L); and a diminishing step size
$\eta_t \to 0$ with $\sum_t \eta_t = \infty$ recovers convergence under (C) alone,
at the cost of the geometric rate. (L) is the price of running an algorithm rather
than a flow.

One constant then pays for three separate things --- the contraction of the mean
block, the uniform peak curvature that makes the widths descend, and the
denominators of the gain matrix in (G) --- which is why it is worth stating
early rather than introducing it where it is most convenient.

\section{What is assumed away}\label{sec:gaps}

Four things, in decreasing order of how much they should worry a reader.

\paragraph{1. Exact gradients.} Everything above assumes the payoff and its
derivatives are available exactly. With sampled estimates the mean block acquires
a noise term that does \emph{not} vanish with the tracking error --- and vanishing
disturbance is exactly the mechanism of \S\ref{sec:vanish} on which the whole
composition rests. The argument would then deliver an $O(\sqrt\eta)$ ball rather
than convergence; a decreasing step size or variance reduction would be needed.
This is the largest gap between the theorem and any implementation.

\paragraph{2. Verifying (C) at scale.} Two quantifiers hide in (C) --- over every
admissible width, and over every reachable opponent policy --- and neither can be
checked by evaluating anything finitely many times. Each is collapsed by a
structural assumption and by nothing else:

\begin{itemize}[leftmargin=1.4em, itemsep=3pt]
\item \emph{The opponent quantifier} is collapsed by a \emph{finite-rank}
  coupling: if $u(a,b) = D_0(a) - D_1(b) + \ip{\varphi(a)}{C\psi(b)}$ then
  $Q_\nu(a) = D_0(a) + \ip{\varphi(a)}{Cy}$ with $y = \E_\nu \psi$, so
  \textbf{the opponent enters through finitely many numbers} --- the averages of
  its features. The quantifier over all probability measures on $\cB$ collapses
  to one over the compact convex \emph{moment body} $\conv \psi(\cB)$, and the
  condition becomes: \emph{check your own opponent-free landscape, then check the
  opponent cannot bend it enough to matter}. Concretely,
  \[
    \sup_{y} \|Cy\| \cdot \sup_{\cA_j}\|\nabla\varphi_s\|
    \;<\; \inf_{r>0} \rho_0(r)
  \]
  suffices, where $\rho_0$ is the modulus of the self-term alone. This is crude
  --- the right-hand side is an infimum over the \emph{whole} cell, attained far
  from the target where the landscape is nearly flat --- so a game can satisfy (C)
  comfortably while failing this test by orders of magnitude. It is a cheap
  sufficient condition, not a characterisation.
\item \emph{The width quantifier} is collapsed, in one dimension and for
  Gaussians, by mode counting: Gaussian smoothing in $1$D cannot \emph{create}
  local maxima, so counting peaks at the floor bounds the count at every larger
  width. \textbf{Above one dimension this fails} --- smoothing can create maxima ---
  and the honest certificate becomes a numerical check across a ladder of widths
  rather than a proof.
\end{itemize}

\noindent
In practice one rarely certifies (C); one tries to \emph{refute} it, and the two
one-line tests of \S\ref{sec:fatal} --- mode count against support size, and the
resolution inequality \eqref{eq:resolution} --- do most of the work.

\paragraph{3. Exploration between cells.} Covering is assumed because the
algorithm cannot achieve it (\S13).

\paragraph{4. Extensive form.} Nothing here transfers to sequential games as a
corollary. The payoff is bilinear in sequence form, but the policy parameterises
behavioural strategies per information set, and the cellwise analysis would have
to be redone with continuation values as the landscape. That is a second theory,
not an extension.

\part{Summary}

\section{The whole thing on one page}

\paragraph{What is proved.} The problem factorises, and the two factors are not
equally hard. \emph{Mixing is free}: the categorical head plays a finite matrix
game exactly, and finite matrix games are solved unconditionally, with no
assumption on the payoff. \emph{Location is everything}: it is local ascent on a
blurred landscape, so it succeeds exactly when each region of the action space has
one attractor (C), it converges in discrete time when the peaks are nondegenerate
(L), the peaks it finds are the equilibrium atoms by (M), it survives the opponent
moving provided curvature dominates coupling, and the four moving blocks compose
provided no round trip amplifies (G). The residual error is representational, and
the annealing mechanism drives it to zero at $\Theta(1/t)$.

\paragraph{What fixes the shape of the theory.} The negative result. A single
spurious attractor produces an exponentially stable non-Nash configuration with an
open basin, immune to the magnet, to extra capacity and to annealing. So no
initialisation-robust class can permit one, and ``exactly $M$ attractors'' is
\emph{forced} rather than chosen.

\paragraph{The practical message.} Failures usually attributed to tuning are
properties of the game, and the cheapest diagnostic is a mode count against a
support size, available before training.

\paragraph{Everything, sorted.}

\begin{center}
\begin{tabular}{@{}p{0.30\textwidth}p{0.14\textwidth}p{0.46\textwidth}@{}}
\hline
what & kind & effect \\
\hline
one attractor per cell (C) & game & required; its failure is fatal and unfixable \\
nondegenerate peaks (L) & game & required for a fixed step; gives the rate \\
peaks are the atoms (M) & game & what makes the limit an equilibrium \\
small gain (G) & game & required; no $\sigma$, no $\eta$ --- checkable in advance \\
curvature $\alpha$ & game & large is better, but small only costs rate \\
coupling strength $c$ & game & capped by $\alpha^2$; sharp peaks tolerate more \\
no self-term (flat game) & game & atoms become unidentifiable; magnet mandatory \\
\hline
covering initialisation & assumed & cannot be achieved by the algorithm \\
exact gradients & assumed & largest gap to implementation \\
\hline
component family & choice & Gaussian buys annealing, costs $O(s^2)$; point mass costs nothing and buys nothing \\
width floor $s^{\min}$ & knob & sets the bias; must be under half the atom spacing; tolerates finer structure \\
capacity $K$ & knob & $K > M$ helps coverage \emph{inside} the class, feeds traps outside it \\
entropy temperature $\tau$ & knob & free: buys stability, does not move the limit \\
Gaussian magnet $\tau^{\mathrm g}$ & knob & pure brake on curved games; mandatory on flat ones \\
step size $\eta$ & knob & needs $\eta < \alpha/L^2$; cancels out of (G) \\
\hline
\end{tabular}
\end{center}

\end{document}
